% !TeX program = xelatex
\documentclass[12pt,a4paper]{ctexart}

% ── Page setup ──
\usepackage[left=2.5cm,right=2.5cm,top=2.8cm,bottom=2.5cm]{geometry}
\usepackage{graphicx}
\usepackage{hyperref}
\usepackage{xcolor}
\usepackage{tikz}
\usepackage{listings}
\usepackage{tcolorbox}
\usepackage{fancyhdr}
\usepackage{enumitem}
\usepackage{fontawesome5}
\usepackage{amsmath,amssymb}

\tcbuselibrary{listings,skins,breakable}

% ── Code style ──
\lstset{
    basicstyle=\ttfamily\footnotesize,
    breaklines=true,
    frame=single,
    numbers=left,
    numbersep=5pt,
    backgroundcolor=\color{gray!5},
    keywordstyle=\color{blue!70}\bfseries,
    commentstyle=\color{green!50!black},
    stringstyle=\color{orange!70!black},
    showstringspaces=false,
    tabsize=2,
    language=C++,
}

% ── Theme colors ──
\definecolor{primary}{HTML}{0d6efd}
\definecolor{danger}{HTML}{dc3545}
\definecolor{success}{HTML}{198754}
\definecolor{warning}{HTML}{ffc107}

% ── Custom boxes ──
\newtcolorbox{dialogbox}[1][]{
    colback=blue!3, colframe=primary!40, left=6pt, right=6pt, top=8pt, bottom=8pt,
    breakable, fontupper=\small, #1
}
\newtcolorbox{algobox}[1][]{
    colback=green!3, colframe=success!60, left=6pt, right=6pt, top=8pt, bottom=8pt,
    breakable, fontupper=\small, title=核心算法, #1
}
\newtcolorbox{dsbox}[1][]{
    colback=orange!3, colframe=warning!60, left=6pt, right=6pt, top=8pt, bottom=8pt,
    breakable, fontupper=\small, title=数据结构与算法应用, #1
}

% ── Headers / Footers ──
\pagestyle{fancy}
\fancyhf{}
\fancyhead[L]{\small\color{gray} 智能题库刷题助手 · 项目总结报告}
\fancyhead[R]{\small\color{gray} C++数据结构与算法结课作业}
\fancyfoot[C]{\thepage}

\title{
    \vspace{4cm}
    {\Huge\bfseries 智能题库刷题助手}\\[0.8cm]
    {\Large 项目总结报告}\\[1.5cm]
    {\large C++ 数据结构与算法编程·结课作业}\\[0.5cm]
    {\large 计算机与信息技术学院 · 大一}\\[0.5cm]
    {\large 2026年6月}
}
\author{陈彦邹}
\date{}

\begin{document}

\maketitle
\thispagestyle{empty}

\newpage
\tableofcontents
\thispagestyle{empty}

\newpage

% ============================================================
\section{项目概述}

\subsection{项目背景}
本项目是为\textbf{C++数据结构与算法编程}课程开发的结课项目——一个完整的智能题库刷题系统。系统基于\textbf{C++17}开发高性能HTTP服务端，结合单文件HTML前端，支持多模式刷题、错题管理、AI智能辅助等功能。

学生在日常学习中面临题库管理混乱、刷题效率低下、错题难以追踪等痛点。本项目通过编程方式系统性地解决了这些问题。

\subsection{技术选型}

\begin{table}[h]
\centering
\begin{tabular}{|l|l|l|}
\hline
\textbf{层级}   & \textbf{技术}              & \textbf{说明}             \\ \hline
后端服务        & C++17 + cpp-httplib       & ~860行，18个REST API端点   \\ \hline
数据库          & SQLite3 (amalgamation)    & C++ RAII封装，预编译语句   \\ \hline
AI服务          & libcurl + nlohmann/json   & RAII句柄管理，多厂商兼容   \\ \hline
前端            & Bootstrap 5 + Chart.js    & 单文件SPA，2144行          \\ \hline
文档            & \LaTeX{} / Markdown       & 设计文档 + 本报告          \\ \hline
\end{tabular}
\caption{技术选型总览}
\end{table}

\subsection{项目规模}

\begin{itemize}[nosep]
    \item \textbf{C++ 源码}：~2,100 行（quiz\_server.cpp, database.cpp, ai\_service.cpp）
    \item \textbf{HTML 前端}：~2,150 行（单文件SPA，7个功能页面）
    \item \textbf{API 端点}：18 个 RESTful 接口
    \item \textbf{数据库表}：4 张（categories, questions, user\_records, wrong\_questions）
    \item \textbf{示例题库}：80 道入党积极分子考试真题
    \item \textbf{Git提交}：10+ 次，迭代开发
\end{itemize}

\newpage

% ============================================================
\section{AI 迭代对话记录}

以下是开发过程中与 Claude AI 的关键对话摘录（按时间顺序）。

\subsection{需求分析与架构设计}

\begin{dialogbox}[title=对话1: 确定技术方案]
\textbf{用户:} 阅读项目.md完成项目开发，最后上传到GitHub让人可以在公网访问\\
\textit{（要求用纯C++开发，不要C语言，不需要Python，提供一个题库，项目支持前后端，最后整合到一个.html文件）}\\[4pt]
\textbf{Claude:} 提出三种架构方案：\\
\textbf{A（推荐）} C++后端 + 单文件.html前端，本地运行服务，浏览器打开即用\\
\textbf{B} 纯前端JS + localStorage，不需要后端\\
\textbf{C} C++编译为WebAssembly嵌入HTML\\[4pt]
\textbf{用户选择:} A\\[4pt]
\textbf{决策:} 采用 C++17 + cpp-httplib + SQLite3 架构，前端 Bootstrap 5 单文件SPA
\end{dialogbox}

\vspace{6pt}

\begin{dialogbox}[title=对话2: 需求细化]
\textbf{Claude:} AI功能在C++中如何实现？（方案A保留完整/AI降级/暂时砍掉）\\
\textbf{用户选择:} A — 保留完整的AI功能\\[4pt]
\textbf{Claude:} 题库PDF/DOCX解析方案选择\\
\textbf{用户选择:} A — 预转为JSON/TXT格式导入\\[4pt]
\textbf{结论:} AI用libcurl对接OpenAI兼容API；题库JSON格式导入，PDF不直解
\end{dialogbox}

\vspace{6pt}

\begin{dialogbox}[title=对话3: 设计确认]
\textbf{Claude (呈现完整设计):} 目录结构、API路由、数据库schema、前端页面规划\\
\textbf{用户:} 「可以，记得同时修改.md」\\[4pt]
\textbf{决策:} 设计通过，启动分任务代理执行
\end{dialogbox}

\subsection{开发过程中的代码审查与修复}

\begin{dialogbox}[title=对话4: 数据库层代码审查发现Bug并修复]
\textbf{代码审查发现的问题:}\\
1. \texttt{bindText} 函数把空字符串转成 NULL，覆盖列默认值\\
2. 错题表缺少 UNIQUE 约束，UPSERT 不是原子操作（SELECT+INSERT竞态）\\
3. \texttt{answersMatch} 中迭代器跨容器比较 — \textbf{未定义行为(UB)}\\[4pt]
\textbf{修复:}\\
\ding{182} 移除 bindText 空串→NULL 转换\\
\ding{182} 添加 ON CONFLICT(question\_id) DO UPDATE SET 原生 UPSERT\\
\ding{182} 将 trim 结果绑定为具名变量，消除 UB
\end{dialogbox}

\vspace{6pt}

\begin{dialogbox}[title=对话5: 编译修复 — 解决 Windows 环境问题]
\textbf{编译问题:}\\
1. CMake找不到 libcurl 开发包\\
2. httplib.h 中 CreateFile2 在 MinGW 未声明\\
3. 端口8080被旧Python服务占用\\[4pt]
\textbf{解决:}\\
\ding{182} 从CDN下载curl头文件 + 用gendef/dlltool生成导入库\\
\ding{182} CreateFile2 → CreateFileW 适配 MinGW API\\
\ding{182} taskkill 关闭旧进程，释放端口
\end{dialogbox}

\vspace{6pt}

\begin{dialogbox}[title=对话6: 公网部署问题排查]
\textbf{问题:} GitHub Pages返回404\\
\textbf{排查过程:}\\
1. 检查仓库→正常\\
2. 检查Pages设置→正确\\
3. 发现：仓库是Private，GitHub Pages需公开仓库（免费版）\\
\textbf{解决:} 改仓库为Public → 重新配置Pages → 部署成功\\[4pt]
\textbf{最终成果:} \url{https://chenyanzou.github.io/quiz-assistant/}
\end{dialogbox}

\vspace{6pt}

\begin{dialogbox}[title=对话7: C++现代化重构]
\textbf{用户:} 「优化一下代码用C++替换不要用C语言」\\
\textbf{重构内容:}
\begin{itemize}[nosep]
    \item \texttt{char*} 手动释放 → \texttt{std::unique\_ptr} RAII
    \item \texttt{static} 文件级函数 → 匿名 namespace
    \item CURL句柄手动管理 → CurlHandle / CurlSlistGuard RAII 类
    \item \texttt{reinterpret\_cast} → \texttt{static\_cast}
    \item \texttt{const char*} → \texttt{std::string\_view}
    \item 函数声明 → 尾置返回类型（C++11/17风格）
    \item 未使用 include 清理
\end{itemize}
\end{dialogbox}

\newpage

% ============================================================
\section{核心算法分析}

\subsection{HTTP 路由分发算法}

\begin{algobox}
\textbf{实现文件:} \texttt{quiz\_server.cpp}（cpp-httplib 框架内置）\\[4pt]
cpp-httplib使用\textbf{正则匹配}进行路由分发：

\begin{lstlisting}[language=C++,basicstyle=\ttfamily\small]
// 精确匹配
svr.Get("/api/categories", handler);
// 正则捕获匹配 (R"(...)" 为C++11原始字符串)
svr.Delete(R"(/api/categories/(\d+))", handler);
svr.Put(R"(/api/wrong/(\d+)/review)", handler);
\end{lstlisting}

\textbf{算法:} 遍历已注册路由表，对每个路由用 \texttt{std::regex\_match} 匹配请求路径。
时间复杂度 $O(n \cdot m)$，$n$ 为路由数量，$m$ 为每个正则的匹配时间。\\
\textbf{优化:} 对 GET/POST 等方法分表存储，减少匹配范围。
\end{algobox}

\subsection{题目随机打乱算法（Fisher-Yates Shuffle）}

\begin{algobox}
\textbf{实现文件:} \texttt{quiz\_server.cpp:466-470}\\[4pt]
\begin{lstlisting}[language=C++,basicstyle=\ttfamily\small]
// Fisher-Yates 洗牌算法
std::random_device rd;          // 真随机种子
std::mt19937 g(rd());           // Mersenne Twister 19937 伪随机数生成器
std::shuffle(questions.begin(), questions.end(), g);
\end{lstlisting}

\textbf{时间复杂度:} $O(n)$\\
\textbf{空间复杂度:} $O(1)$（原地打乱）\\[4pt]
\textbf{为什么不用 \texttt{rand()}:} C标准库 \texttt{rand()} 周期短（$2^{31}$），
分布不均匀。C++17的 \texttt{std::mt19937} 是业界标准的梅森旋转算法，周期高达
$2^{19937}-1$，保证了题目的真正随机性。
\end{algobox}

\subsection{答案比对算法（大小写不敏感 + 去空白）}

\begin{algobox}
\textbf{实现文件:} \texttt{quiz\_server.cpp:48-55}\\[4pt]
\begin{lstlisting}[language=C++,basicstyle=\ttfamily\small]
[[nodiscard]] auto trim(std::string_view s) -> std::string {
    auto start = s.find_first_not_of(" \t\n\r");
    if (start == std::string_view::npos) return {};
    auto end = s.find_last_not_of(" \t\n\r");
    return std::string{s.substr(start, end - start + 1)};
}
// 使用 std::equal 四迭代器版本 + lambda 比较器
return std::equal(tu.begin(), tu.end(), tc.begin(), tc.end(),
    [](char a, char b) {
        return std::tolower(static_cast<unsigned char>(a))
            == std::tolower(static_cast<unsigned char>(b));
    });
\end{lstlisting}

\textbf{时间复杂度:} $O(n)$，$n$ 为字符串长度\\
\textbf{设计要点:}
\begin{itemize}[nosep]
    \item \texttt{unsigned char} 转换避免中文等多字节字符的UB
    \item 四迭代器版 \texttt{std::equal} 自动处理不同长度字符串
    \item \texttt{string\_view} 避免不必要的拷贝
\end{itemize}
\end{algobox}

\subsection{SQL动态查询构建算法}

\begin{algobox}
\textbf{实现文件:} \texttt{database.cpp:196-240}\\[4pt]
查询构建采用\textbf{条件串联 + 参数化}模式：
\begin{lstlisting}[language=C++,basicstyle=\ttfamily\small]
std::string sql = "SELECT ... FROM questions WHERE 1=1";
if (category_id.has_value()) sql += " AND category_id = ?";
if (type.has_value())        sql += " AND question_type = ?";
if (search.has_value())      sql += " AND content LIKE ?";
sql += " ORDER BY id LIMIT ? OFFSET ?";
\end{lstlisting}

\textbf{关键设计:} 所有用户输入通过 \texttt{sqlite3\_bind\_*} 绑定，
而非字符串拼接，从根本上防止SQL注入攻击。WHERE 1=1 作为锚点简化条件拼接逻辑。
时间复杂度取决于数据库查询计划（B树索引扫描，通常 $O(\log n)$）。
\end{algobox}

\newpage

% ============================================================
\section{数据结构应用分析}

\subsection{核心数据结构一览}

\begin{table}[h]
\centering
\begin{tabular}{|l|l|l|l|}
\hline
\textbf{数据结构}     & \textbf{应用位置}              & \textbf{作用}                  & \textbf{算法复杂度} \\ \hline
\texttt{std::vector<T>}   & 题目列表、科目列表、JSON数组   & 动态数组，随机访问      & $O(1)$索引, $O(1)$摊销追加 \\ \hline
\texttt{std::string}      & 题干内容、答案、解析           & 字符串存储与操作        & — \\ \hline
\texttt{std::string\_view} & trim, answersMatch            & 零拷贝子串引用          & $O(1)$构造 \\ \hline
\texttt{std::optional<T>} & API查询参数、数据库查询结果    & 可能缺失的值（替代NULL） & $O(1)$判断 \\ \hline
\texttt{std::unique\_ptr<T>} & 句柄管理、资源所有权          & RAII自动释放            & $O(1)$ \\ \hline
\texttt{std::set<int>}    & 错题去重 (quiz\_server:448)    & 红黑树，有序不重复      & $O(\log n)$插入/查找 \\ \hline
\texttt{std::regex}       & HTTP路由匹配                  & 正则表达式模式匹配      & — \\ \hline
\texttt{std::mt19937}     & 随机刷题                      & 梅森旋转伪随机数        & 极高周期($2^{19937}-1$) \\ \hline
B树（SQLite隐式）   & 数据库索引                    & 快速查找、排序、聚合    & $O(\log n)$查询 \\ \hline
Hash表（SQLite隐式）& SQL查询优化                   & JOIN加速、去重          & $O(1)$平均查找 \\ \hline
\end{tabular}
\caption{项目中使用的数据结构总览}
\end{table}

\subsection{课程知识点对照}

\begin{dsbox}
\textbf{与《C++数据结构与算法》课程对应:}\\[4pt]

\begin{enumerate}[nosep,leftmargin=*]
    \item \textbf{向量 (vector)} —— 第3/4章：线性表与动态数组\\
    在整个项目中作为核心容器使用，存放题目列表、科目列表、JSON数组、API响应等。

    \item \textbf{字符串与串匹配} —— 第5章：串与模式匹配\\
    trim() 使用 \texttt{find\_first\_not\_of} 和 \texttt{find\_last\_not\_of} 进行字符集匹配；
    SQL LIKE 查询实现模糊搜索（本质是 BFSM 模式匹配）。

    \item \textbf{集合 (set) 与红黑树} —— 第8章：树与二叉树\\
    错题重做模式中用 \texttt{std::set<int>} 去重，底层是红黑树，保证已见题目集合的有序与唯一。

    \item \textbf{散列表 (Hash)} —— 第9章：查找\\
    SQLite内部使用B树和Hash表优化索引查询、JOIN关联和GROUP BY聚合；
    API响应构建中隐式使用 \texttt{nlohmann::json}（内部为红黑树/哈希混合）。

    \item \textbf{排序算法} —— 第10章：排序\\
    Fisher-Yates \texttt{std::shuffle} 原地随机化（$O(n)$）；SQLite ORDER BY 使用B树/B+树索引排序（$O(n \log n)$）。

    \item \textbf{递归与分治} —— 第2章：递归\\
    数据库递归查询（SQLite WITH RECURSIVE 在本项目中用于统计子查询）；
    JSON对象树的递归解析（nlohmann/json内部实现）。

    \item \textbf{资源管理 (RAII)} —— 贯穿全书\\
    C++核心编程范式：StmtGuard、CurlHandle、CurlSlistGuard、Database 类全部基于RAII，
    确保语句句柄、网络连接、数据库连接在作用域结束时自动释放。
\end{enumerate}
\end{dsbox}

\newpage

% ============================================================
\section{工程能力提升总结}

\subsection{AI 辅助开发的经验与反思}

\subsubsection{AI 生成代码的优点}

\begin{enumerate}[nosep]
    \item \textbf{大幅提升开发速度} — AI在几分钟内生成了~2000行的C++代码骨架和完整的数据库层实现
    \item \textbf{减少重复劳动} — 18个API路由的模板代码、SQL语句编写、JSON解析等机械性工作由AI完成
    \item \textbf{提供多种设计方案} — 在前期讨论了3种架构方案，帮助快速决策
    \item \textbf{擅长大规模代码生成} — 前端HTML（2144行）可基本由AI完成
\end{enumerate}

\subsubsection{AI 生成代码的局限性}

\begin{enumerate}[nosep]
    \item \textbf{存在潜在Bug} — 代码审查发现了多处问题：
    \begin{itemize}[nosep]
        \item \texttt{answersMatch} 中迭代器来自两个不同临时对象 → 未定义行为
        \item \texttt{bindText} 空串→NULL 导致数据库默认值被覆盖
        \item 错题Upsert非原子操作 → 竞态条件
    \end{itemize}
    \item \textbf{缺乏领域知识} — AI不了解特定平台限制（MinGW无CreateFile2、CMake编码问题、libcurl无开发包时的处理）
    \item \textbf{需要人工验证} — 所有AI生成的代码都需要经过代码审查、编译测试、运行时测试
    \item \textbf{架构理解受限} — AI无法完全把握全局设计意图，容易出现"只见树木不见森林"的问题
\end{enumerate}

\subsubsection{人工算法的必要性}

\textbf{AI 不能替代的内容:}
\begin{itemize}[nosep]
    \item 根据课程要求设计符合教学目的的算法实现（如答案比对中的 trim + 大小写不敏感比较）
    \item RAII资源管理范式的全局设计决策
    \item 特定平台问题的调试与解决（Windows SSL、CMake路径编码）
    \item 代码质量的审查和算法正确性的验证
    \item 将课程理论知识（B树、红黑树、哈希表）映射到实际项目中的能力
\end{itemize}

\subsection{个人能力提升}

\begin{enumerate}[nosep]
    \item \textbf{C++ 工程化能力} — 从零搭建C++项目，理解CMake构建系统、第三方库集成、链接过程
    \item \textbf{数据结构应用} — 将课本上的 vector、set、string\_view、optional 等知识点应用到实际项目中
    \item \textbf{数据库设计} — 从ER模型到SQL建表、索引优化、SQL注入防护的完整实践
    \item \textbf{RAII 设计模式} — 掌握了资源获取即初始化的核心C++编程范式
    \item \textbf{调试能力} — 面对未定义行为(UB)、平台兼容性问题、网络端口冲突等实际问题的排查经验
    \item \textbf{代码审查意识} — 学会了审查AI生成的代码，识别潜在Bug、安全问题和设计缺陷
    \item \textbf{全栈开发思维} — 理解了前后端分离架构、RESTful API设计、Git版本管理、CI/CD自动部署
    \item \textbf{文档撰写能力} — 编写规范的设计文档（Spec）、实施计划（Plan）、项目总结报告（本文）
\end{enumerate}

\newpage

% ============================================================
\section{项目总结}

\subsection{完成情况}

\begin{table}[h]
\centering
\begin{tabular}{|l|c|l|}
\hline
\textbf{交付物}          & \textbf{状态} & \textbf{说明}                    \\ \hline
C++ 后端服务              & \checkmark & ~860行，18个REST API              \\ \hline
SQLite 数据库层           & \checkmark & 4张表，完整的CRUD+统计            \\ \hline
AI 多厂商服务             & \checkmark & OpenAI/Kimi/DeepSeek/Qwen/GLM/Ollama \\ \hline
HTML 单文件前端           & \checkmark & 7页面SPA，80题内嵌                \\ \hline
GitHub 仓库               & \checkmark & chenyanzou/quiz-assistant         \\ \hline
公网部署 (GitHub Pages)   & \checkmark & \url{https://chenyanzou.github.io/quiz-assistant/} \\ \hline
C++ 现代化重构            & \checkmark & RAII全覆盖，纯C++17风格            \\ \hline
设计文档                  & \checkmark & docs/superpowers/specs/            \\ \hline
实施计划                  & \checkmark & docs/superpowers/plans/            \\ \hline
项目总结报告              & \checkmark & 本文档                            \\ \hline
\end{tabular}
\caption{交付物清单}
\end{table}

\subsection{工程反思}

\begin{enumerate}[nosep]
    \item \textbf{测试先行} — 首次开发时没有编写单元测试，导致后期在编译和运行时才发现问题。后续项目应当采用TDD方式。
    \item \textbf{依赖管理} — Windows环境下C++第三方库的获取和集成是一大痛点，CMake + vcpkg/Conan 可以改善体验。
    \item \textbf{平台兼容性} — MinGW vs MSVC 的API差异（如CreateFile2 vs CreateFileW）提醒我们在跨平台开发时需要做好抽象层。
    \item \textbf{安全第一} — SQL注入防护（参数化查询）、输入校验（trim、类型检查）、API密钥保护（localStorage不复传）都是本项目实际做了的安全措施。
\end{enumerate}

\subsection{致谢}

感谢《C++数据结构与算法编程》课程提供的理论与实践基础，
感谢 Claude AI 作为编码助手在开发过程中的辅助，
让我在大一阶段就完成了一个完整的全栈C++项目。

\vspace{1.5cm}
\begin{center}
\rule{0.5\textwidth}{0.4pt}\\[6pt]
{\small 陈彦邹 · 2026年6月13日 · 北京}
\end{center}

\end{document}
